\documentclass[aps,prx,twocolumn,superscriptaddress,nofootinbib,floatfix]{revtex4-2}

\usepackage{amsmath,amssymb}
\usepackage{booktabs}
\usepackage{xcolor}
\usepackage{hyperref}   % loaded last, after xcolor
\hypersetup{colorlinks=true,linkcolor=blue,citecolor=blue,urlcolor=blue}

\begin{document}

\title{A pre-registered entanglement-forged H$_4$ benchmark on trapped-ion
hardware: accurate noise modelling, incomplete error correction}

\author{Venkata Rao Allu}
\email{alluvenkat11@gmail.com}
\affiliation{Independent researcher}

\date{\today}

\begin{abstract}
We report a pre-registered benchmark of an entanglement-forged H$_4$
ground-state calculation on IonQ Forte Enterprise~1, a live trapped-ion
processor. The analysis was frozen and publicly committed before any
hardware data existed; the run was a single submission of $81$ circuits at
$1100$ shots with device debiasing explicitly disabled. Two findings are in
tension and we report both. A three-parameter analytic noise model, fixed
in advance, predicted the \emph{uncorrected} energy error to within $1\%$
($23.84$ against $24.06$~kcal/mol) and the postselection retention to
$0.4\%$. The correction built on that same model then failed: the
oracle-free estimator returned $3.46$~kcal/mol against $0.83$ predicted, a
factor of $4.16$, missing chemical accuracy by $3.5\times$ despite removing
$85\%$ of the raw error. Goodness of fit degraded from
$\chi^2/\mathrm{dof}\approx2$ on every noise model to $5.93$ on hardware,
in all $20$ bootstrap resamples, so the failure was detectable from the
hardware counts alone without reference to the known answer. Decomposing
all $732$ residuals by measurement group and Pauli weight localises the
discrepancy to high-weight $XY$ strings carrying the largest corrections
and to a signed offset in one group. Refitting a single correction-magnitude
parameter, blind to the reference energy, shows the model
\emph{under}-corrects by $\approx1.8\times$; repairing this removes only a
third of the excess and makes the energy slightly worse, so fit quality and
accuracy are decoupled. We conclude that a noise model calibrated on raw
observables must not be assumed predictive of corrected ones, that
bootstrap resampling understates submission-to-submission variance, and
that at $1100$ shots the oracle-free estimator is shot-noise limited rather
than noise limited, so mitigation alone cannot close the gap.
\end{abstract}

\maketitle

\section{Introduction}

Variational quantum eigensolvers~\cite{peruzzo2014,mcclean2016} are the
most-studied near-term application of quantum processors, and the usual
measure of success is whether a computed molecular energy falls within
chemical accuracy, conventionally $1$~kcal/mol
($1.6\times10^{-3}$~Hartree), of a classical reference. A large
literature reports such agreement, most often on simulators with device
noise models, sometimes on hardware with substantial error mitigation.

This paper takes a deliberately narrow and unglamorous question. Given a
noise model that has been calibrated until it reproduces a device's
\emph{raw}, unmitigated behaviour, does the error correction built on that
model transfer to the device? We answer it for one molecule, one device,
one submission, with the analysis fixed in advance, and we report the
answer whichever way it came out.

It came out negative in a specific and, we argue, informative way. The
noise model's prediction of the raw error was nearly exact. Its prediction
of the corrected error was wrong by more than a factor of four. The
goodness-of-fit statistic flagged the failure --- it tripled relative to
every simulator rehearsal --- which means the failure was detectable from
the hardware data alone, without reference to the known answer. We regard
that last point as the most portable result here: a mitigation pipeline
should report a calibrated fit statistic, because that statistic
distinguishes ``the correction worked'' from ``the correction produced a
number'' without consulting an oracle.

Three practices made this possible and we recommend them. The analysis
code was committed before the data existed, so no post hoc choice could
be made. Every estimator was reported, including the one that failed.
And the goodness-of-fit was computed and published alongside each energy.

\section{The electronic-structure engine}

All classical quantities were computed from a from-scratch implementation:
contracted Gaussian basis evaluation, one- and two-electron integrals via
the Obara--Saika recursions, restricted Hartree--Fock, integral
transformation to the molecular-orbital basis, and the Jordan--Wigner
mapping to a qubit Hamiltonian. Nothing is taken from a quantum-chemistry
package at runtime. Correctness is established by comparison against
PySCF~\cite{pyscf} on the same geometries and basis, and the engine
reproduces its Hartree--Fock and full configuration interaction (FCI)
energies to numerical precision. We emphasise this only because it
determines what the error bars mean: the reference energy against which
all errors below are quoted is an exact FCI diagonalisation in the same
basis, computed by the same code, not an external literature value.

The system is H$_4$ in a linear chain at $1.0$~\AA{} spacing, STO-3G,
$4$ electrons in $4$ spatial orbitals, giving an $8$-qubit Hamiltonian.

\section{Methods}

\subsection{Entanglement forging}

The $8$-qubit problem is decomposed by entanglement
forging~\cite{eddins2022}, which writes the ground state as a Schmidt
decomposition across a bipartition into two $4$-qubit registers,
\begin{equation}
|\psi\rangle = \sum_{n=1}^{K} \lambda_n\, |u_n\rangle_\alpha \otimes
|v_n\rangle_\beta ,
\end{equation}
so that the energy becomes a weighted sum of expectation values measured
on one $4$-qubit register at a time, over Schmidt basis states
$|u_n\rangle$ and their pairwise superpositions
$(|u_n\rangle+|u_m\rangle)/\sqrt2$. For H$_4$ at this geometry the Schmidt
rank is effectively $K=6$, giving $21$ distinct preparation ``slots''
($6$ singles and $15$ pairs).

This is the first of several places where exact information enters the
construction, and we flag each one explicitly. The state-preparation
angles for each slot are obtained by classically fitting the circuit to
reproduce the exact FCI Schmidt vectors, and the Schmidt coefficients
$\lambda_n$ used to assemble the energy are the exact ones. The
calculation therefore measures whether hardware can reproduce a state
whose identity is already known; it is not a blind variational search.
Section~\ref{sec:limits} states the consequences.

\subsection{Native-gate compilation}

Circuits are compiled directly to IonQ's native gate set --- $\mathrm{GPi}$,
$\mathrm{GPi2}$, and the two-qubit $\mathrm{ZZ}$ interaction --- rather than
to a CNOT-based abstraction that the provider then retranslates. Angles are
expressed in turns, with the $\mathrm{ZZ}$ angle restricted to
$[0,0.25]$. Where a negative interaction angle is required we use the
identity $X\!\cdot\!\mathrm{ZZ}(\theta)\!\cdot\!X = \mathrm{ZZ}(-\theta)$
with $X=\mathrm{GPi}(0)$, since the inverse method of the provider's gate
objects does not negate the parameter.

\subsection{General-commuting measurement grouping}

The $\alpha$-register Hamiltonian requires $273$ Pauli expectation values.
Grouping them into sets of mutually commuting --- not merely
qubit-wise commuting --- operators and diagonalising each set with a
Clifford circuit reduces this to $4$ measurement groups. Across $21$ slots
this gives $84$ circuits rather than $273$, a $3.25\times$ reduction. The
grouping is obtained by exact graph colouring and verified to
$6.7\times10^{-16}$ against direct evaluation before use.

\subsection{Ancilla-parity leakage detection}

One ancilla qubit records the total parity of the $4$-qubit register. The
physical state lives in a fixed particle-number sector, so shots whose
ancilla reads $1$ have provably left that sector and are discarded. The
retained fraction is reported for every run and is itself a prediction the
noise model makes.

\subsection{Analytic group-aware correction}

The correction below is a forward analytic inverse of an assumed noise
channel, in the spirit of zero-noise extrapolation and probabilistic error
cancellation~\cite{temme2017} and of the sparse Pauli--Lindblad models used
to learn such channels on hardware~\cite{vandenberg2023}, but with the
channel parameters fixed from independent characterisation rather than
learned on the target circuits.

For each slot and measurement group we compute, from the intended circuit
and a three-parameter noise model, two quantities $A_\ell$ and $B_\ell$
per Pauli label $\ell$: the expectation the noisy circuit would produce
and the expectation the ideal circuit would produce. The measured value is
then corrected multiplicatively,
\begin{equation}
\langle P_\ell \rangle_{\mathrm{corr}}
= \langle P_\ell \rangle_{\mathrm{raw}} \cdot \frac{B_\ell}{A_\ell}
\equiv \langle P_\ell \rangle_{\mathrm{raw}} \cdot r_\ell .
\label{eq:corr}
\end{equation}
The three noise parameters --- a $\mathrm{ZZ}$ depolarising rate
$p_{ZZ}=0.014593$, a $\mathrm{GPi}$ rate $p_{\mathrm{GPi}}=0.000119$, and a
$\mathrm{GPi2}$ rate $p_{\mathrm{GPi2}}=0.0004$ --- were fixed before the
hardware run from independent characterisation and never retuned on
hardware data. The correction ratios $r_\ell$ for this experiment lie
between $1.00$ and about $1.20$.

\subsection{Three estimators}

From the same counts we compute three energies, and we report all three.

\emph{Raw}: a per-slot physical-state fit to the postselected but
uncorrected expectation values.

\emph{Frame-free}: the same per-slot fit applied to the corrected values
of Eq.~\eqref{eq:corr}. No exact Schmidt vector enters the
reconstruction. This is the oracle-free estimator and we treat it as the
headline number.

\emph{Shared-frame}: a single $K\times K$ orthogonal frame is fit jointly
across all $21$ slots, initialised at the exact frame, with the
reconstruction basis spanned by the exact FCI Schmidt vectors. This
estimator pools information across slots and is substantially tighter,
but it depends on the exact state in two additional ways beyond the
circuit construction. We report it because it was pre-registered, and we
do not use it to claim chemical accuracy.

\subsection{Goodness of fit}

For the shared-frame fit we compute a shot-noise-weighted residual for
every (slot, label) pair,
\begin{equation}
R_{s\ell} = \frac{\langle \hat v_s | P_\ell | \hat v_s\rangle
- \langle P_\ell\rangle_{\mathrm{corr}}}{\sigma_{s\ell}},
\qquad
\sigma_{s\ell} = |r_\ell|\sqrt{\frac{1-\langle P_\ell\rangle^2_{\mathrm{raw}}}{N_{s\ell}}},
\end{equation}
with $N_{s\ell}$ the retained shot count, and report
$\chi^2/\mathrm{dof} = \sum R^2 / (n - N_\theta)$. Labels whose correction
drives the value to near zero ($|r_\ell| < 0.01$) are excluded and counted.
A well-specified model with correct error bars gives
$\chi^2/\mathrm{dof}\approx1$.

\section{Pre-registration and run conditions}

The analysis script was committed before the hardware job was submitted
and run unmodified on the returned counts. The conditions below were
agreed with IonQ in advance and are stated because two of them materially
affect the result.

\begin{itemize}
\item \textbf{Device debiasing disabled}, passed explicitly per circuit.
With debiasing enabled the same job is priced at approximately $\$13{,}600$
rather than $\$3{,}882$. Debiasing also averages over qubit permutations,
which would confound a correction model calibrated to specific gates.
\item \textbf{$1100$ shots per circuit}, fixed in advance.
\item \textbf{A single submission}, with no resubmission permitted. No
circuit was re-run and no outlier was discarded.
\item $81$ new circuits, plus $3$ reused from an earlier job on the same
device at $2000$ shots, verified to rebuild with identical gate counts.
\end{itemize}

The job was submitted 2026-10-14 05:34~UTC and completed 2026-10-16
03:36~UTC. Cost was $\$3{,}882.21$ by the provider's estimate endpoint and
$\$3{,}875.18$ by an independent rate-card calculation, a $\$7.03$
discrepancy. Single-qubit gates accounted for $58.9\%$ of the total
($12{,}582$ one-qubit and $1{,}302$ two-qubit gates), which is worth
noting because circuit-optimisation effort in this architecture is often
directed at two-qubit count alone.

Two pre-registered gates were checked before analysis: the Schmidt sign
convention reproduced stored ideal data to $3.2\times10^{-5}$ (limit
$10^{-3}$), and all $84$ required circuits were present.

\section{Results}

\subsection{Hardware against noise models}

Table~\ref{tab:main} gives the three estimators on hardware and on three
noise models run through the identical pipeline. We stress that
\emph{aria-1} and \emph{forte-1} denote device noise \emph{models} offered
for simulation, not hardware; the Aria-1 device has been retired and no
submission to it is possible. Only the final row is hardware.

\begin{table}[t]
\caption{Energy error against exact FCI, in kcal/mol, and shared-frame
goodness of fit. Chemical accuracy is $1.0$. Only the last row is a
physical device.}
\label{tab:main}
\begin{ruledtabular}
\begin{tabular}{lcccc}
 & Raw & Frame-free & Shared-frame & $\chi^2/\mathrm{dof}$ \\
\colrule
ideal model        & $1.96$  & $0.048$ & $0.117$ & $\approx1$ \\
aria-1 model       & $22.53$ & $1.041$ & $0.707$ & $\approx2$ \\
forte-1 model      & $24.06$ & $0.832$ & $0.394$ & $2.68$ \\
\textbf{Forte Enterprise 1} & $\mathbf{23.84}$ & $\mathbf{3.46}$ &
$\mathbf{0.86}$ & $\mathbf{5.93}$ \\
\end{tabular}
\end{ruledtabular}
\end{table}

The raw agreement is the strongest quantitative result in this work. The
hardware returned $23.84$~kcal/mol where the forte-1 model predicted
$24.06$, a ratio of $0.99$. Postselection retention was $0.894$ against
$0.898$ predicted. A three-parameter model, fixed in advance, reproduced
both the uncorrected error of an $8$-qubit forged calculation and its
leakage rate on a live device.

The corrected agreement fails. The frame-free estimator returned
$3.46$~kcal/mol where the same model predicted $0.832$, a factor of
$4.16$. The correction is not useless --- it removes $85\%$ of the raw
error, $23.84\to3.46$ --- but it does not reach chemical accuracy, missing
by $3.5\times$.

\subsection{Uncertainty}

Uncertainties come from a nonparametric bootstrap in which every circuit's
counts are resampled at its own retained shot total and the entire
pipeline, including the frame fit, is re-run ($N=20$):
\begin{align}
\text{raw} &: 19.49 \pm 5.65, \quad [14.48,\,28.90]_{5-95} \nonumber\\
\text{frame-free} &: 2.92 \pm 1.90, \quad [0.39,\,6.47]_{5-95} \nonumber\\
\text{shared-frame} &: 1.055 \pm 0.419, \quad [0.42,\,1.87]_{5-95}
\nonumber
\end{align}
The shared-frame point estimate of $0.864$ lies below $1$~kcal/mol, but
its bootstrap distribution is centred \emph{above} the threshold and $13$
of $20$ resamples exceed it. We therefore state that chemical accuracy is
\emph{consistent with} this dataset on the shared-frame estimator, and is
not established by it. The frame-free estimator exceeds $1$~kcal/mol in
$16$ of $20$ resamples.

\subsection{The fit statistic detected the failure}
\label{sec:chi2}

$\chi^2/\mathrm{dof}$ was $0.9$--$1.2$ on the ideal model, $1.6$--$2.7$ on
the two noisy models, and $5.93$ on hardware, with every one of the $20$
bootstrap resamples in $[6.41,7.22]$. The excess is therefore a property
of the data, not of one fit. Because this statistic requires no reference
energy, the inadequacy of the correction model was visible from the
hardware counts alone. Any pipeline reporting a mitigated energy can
report this number, and we suggest that one which does not should be
treated as unaudited.

\section{Localising the discrepancy}

\subsection{Residual decomposition}

We re-derived the identical $732$-term residual vector and tagged each
term with its slot, measurement group, Pauli label and Pauli weight.
Table~\ref{tab:resid} gives mean $R^2$ by group and weight. (The design
measures each Pauli label in exactly one group, so cells are
structurally empty rather than missing.)

\begin{table}[t]
\caption{Mean squared residual $\langle R^2\rangle$ by measurement group
and Pauli weight. Uniform values near $1$ would indicate a well-specified
model.}
\label{tab:resid}
\begin{ruledtabular}
\begin{tabular}{lcccc}
 & $w{=}1$ & $w{=}2$ & $w{=}3$ & $w{=}4$ \\
\colrule
group 0 & $3.59$ & $\mathbf{10.14}$ & --- & --- \\
group 1 & --- & $2.80$ & --- & $4.70$ \\
group 2 & --- & --- & $3.98$ & $\mathbf{13.58}$ \\
group 3 & --- & $3.35$ & --- & $4.77$ \\
\end{tabular}
\end{ruledtabular}
\end{table}

Two structures stand out.

\emph{The residual grows with the magnitude of the correction applied.}
Binning by correction ratio $r_\ell$: labels with
$r_\ell\in[1.00,1.05)$ give $\langle R^2\rangle = 3.69$, while labels with
$r_\ell\in[1.15,1.25)$ give $8.51$ and carry $48\%$ of the total $\chi^2$
from $33\%$ of the terms. The twelve largest individual residuals are all
weight-$4$ $XY$-type strings ($YYXX$, $YXXY$, $XYYX$, $XXYY$, $YYYY$,
$XXXX$), each with $r_\ell\approx1.19$.

\emph{Group $0$ carries a signed offset rather than scatter.} Its mean
\emph{signed} residual is $-1.32$ at weight $1$ and $-1.73$ at weight $2$.
Stochastic misestimation averages to zero; a systematic one does not.

\subsection{Testing the magnitude hypothesis}

The first structure suggests a simple repair: that
Eq.~\eqref{eq:corr} overshoots, so that
\begin{equation}
\langle P_\ell\rangle_{\mathrm{corr}}(\alpha)
= \langle P_\ell\rangle_{\mathrm{raw}}\,\bigl[1 + \alpha\,(r_\ell-1)\bigr]
\label{eq:alpha}
\end{equation}
with $\alpha<1$ would reduce $\chi^2$. Here $\alpha=1$ recovers the
pre-registered pipeline and $\alpha=0$ applies no correction. Crucially,
$\alpha$ is selected by minimising $\chi^2/\mathrm{dof}$ alone; the
reference energy is never an input, preserving the blindness of every
correction-parameter choice in this work. Because fitting $\alpha$ on the
same residuals whose $\chi^2$ is then quoted would reduce it by
construction, we also fit $\alpha$ on half the slots and evaluate
$\chi^2$ on the held-out half.

Two technical points must be settled before the scan is meaningful.

\emph{The exclusion rule is $\alpha$-dependent.} The pre-registered
pipeline discards a label when $|\langle P\rangle_{\mathrm{corr}} /
\langle P\rangle_{\mathrm{raw}}| < 0.01$. Under Eq.~\eqref{eq:alpha} that
quantity is $1+\alpha(r_\ell-1)$, which for the $24$ labels with
$r_\ell\approx0$ is small \emph{only} when $\alpha\approx1$. Recomputing
the rule per $\alpha$ therefore excludes $24$ labels at $\alpha=1$ and
none elsewhere, so $\chi^2$ at different $\alpha$ would be evaluated over
different residual sets, and the pre-registered value would enjoy an
exclusion unavailable to its comparators. We freeze the mask at its
$\alpha=1$ membership, so all values below are computed on the same $732$
terms.

\emph{The $\chi^2$ objective has poles in $\alpha$.} The pre-registered
pipeline propagates shot noise through the correction as
$\sigma_{s\ell}\propto|1+\alpha(r_\ell-1)|$, which is correct for one given
correction but vanishes at $\alpha = 1/(1-r_\ell)$. In this dataset $26$
labels have such a pole within $\alpha\in(0,3]$: $24$ at $\alpha=1.0$ --- the
masked set above --- and $2$ at $\alpha=1.2$, which accounts for an
otherwise unexplained $\chi^2$ spike at that point. The shot noise of a raw
measurement does not change because the correction hypothesis changed, so we
hold $\sigma$ fixed at its $\alpha=1$ value throughout the scan. With both
controls the curve is smooth and monotone about its minimum.

\emph{The fit is converged.} Holding the data fixed at $\alpha=1$ and
varying only the optimiser over $5$ restart counts ($4$ to $64$) and $5$
random seeds, all $25$ configurations returned
$\chi^2/\mathrm{dof}=5.929$ and $0.864$~kcal/mol, with total spreads of
$3.7\times10^{-6}$ and $7\times10^{-6}$~kcal/mol respectively. The pre-registered numbers are not an
artifact of a particular restart draw. The same invariance shows the
random restarts are inert: the fit is initialised at the exact frame and
converges there, which is a further exact-state dependence.

With those controls in place, $\chi^2$ falls as $\alpha$ \emph{increases}:

\begin{center}
\begin{tabular}{lccc}
\toprule
 & $\alpha$ & $\chi^2/\mathrm{dof}$ & err (kcal/mol) \\
\midrule
no correction      & $0.0$ & $14.84$ & $1.017$ \\
pre-registered     & $1.0$ & $5.93$  & $0.864$ \\
blind optimum      & $1.8$ & $3.97$  & $0.940$ \\
over-corrected     & $3.0$ & $7.64$  & $1.166$ \\
\bottomrule
\end{tabular}
\end{center}

The hypothesis we set out to test --- that the correction overshoots ---
is therefore refuted in its direction. The correction \emph{under}-corrects
by roughly a factor of $1.8$. Held-out validation confirms this is not
overfitting: fitting $\alpha$ on half the slots and scoring the other half
improves $\chi^2/\mathrm{dof}$ from $4.97$ to $3.89$ in one fold
($\alpha=1.9$) and from $7.04$ to $4.41$ in the other ($\alpha=1.7$).

Allowing a separate $\alpha_g$ per measurement group improves the
in-sample fit further, to $\chi^2/\mathrm{dof}=3.75$ at
$\alpha_g=(1.9,1.4,2.0,1.3)$, and yields the lowest energy we obtain from
this data, $0.637$~kcal/mol. We do not advance that number. Under the
earlier, $\alpha$-dependent weighting the same model moved to
$(2.0,1.5,2.1,1.6)$ and $0.678$~kcal/mol, and its held-out performance was
no better than the global model's in one fold and worse in the other, with
fitted vectors swinging between $(2.0,1.7,2.4,1.9)$ and
$(2.1,1.3,1.9,1.1)$. Four parameters are not determined by this dataset,
and the group structure is not stable enough to report as a result. The supported statement is the single global
rescaling; the group structure is suggestive only. We note that the two
groups demanding the largest correction, groups $0$ and $2$, are the two
the residual decomposition independently flagged as worst.

Three conclusions follow, and the third is unwelcome. The correction
model's magnitude is genuinely miscalibrated, in the opposite direction to
our hypothesis. Repairing it removes only about a third of the $\chi^2$
excess, $5.93\to3.97$ against a target near unity, so the misspecification
extends beyond a scalar factor. And the better-fitting correction returns a
\emph{worse} energy, $0.940$ against $0.864$. Fit
quality and accuracy are decoupled here, which bounds the claim of
Sec.~\ref{sec:chi2}: $\chi^2$ is a reliable detector of model
misspecification, and not a proxy for accuracy.

A final observation constrains how much of this matters. At $\alpha=0$,
with no correction applied at all, the shared-frame estimator returns
$1.017$~kcal/mol against $0.864$ with the full correction. The entire
analytic correction stack is worth $0.15$~kcal/mol on that estimator; the
joint frame fit supplies the rest. Since the frame fit is the component
carrying the exact-state dependence, the shared-frame result rests
primarily on oracle information rather than on the physics-based
correction. The frame-free estimator, where the correction does the work
unaided, is the one that missed by $4.16\times$.

\section{Discussion}

\subsection{Raw accuracy does not imply corrected accuracy}

The central methodological claim of this paper is narrow and, we believe,
well supported by the single comparison above: a noise model validated
against raw device behaviour was wrong by $4.16\times$ once its own
correction was applied. The raw expectation value is a smooth, heavily
averaged functional of the noise channel; the correction is a ratio of two
such functionals, and ratios amplify misspecification. A noise model's
agreement on uncorrected observables should therefore be treated as a
necessary and distinctly insufficient condition for trusting its
mitigation. Where a projection is needed, the corrected-observable
prediction should be read as an optimistic bound.

\subsection{Bootstrap understates submission variance}

Before the hardware run, bootstrap resampling of earlier data predicted a
frame-free maximum of $0.90$~kcal/mol and a $100\%$ pass rate at
$1$~kcal/mol for this design. Fresh, independent simulator submissions of
the identical design at the same shot count returned up to
$2.46$~kcal/mol. Bootstrap resamples a single submission and therefore
captures shot noise but not drift between submissions; the latter was the
larger term. Reported bootstrap intervals in this setting should be
regarded as lower bounds on the true spread.

\subsection{A shot-noise floor, not a noise floor}

The most consequential constraint we encountered is not device error. On
one fresh $1100$-shot submission the \emph{zero-noise} ideal simulator
itself returned $1.225$~kcal/mol on the frame-free estimator. With no
device error of any kind, the estimator missed chemical accuracy from
counting statistics alone. Improved mitigation cannot repair this.
Scaling as $N^{-1/2}$ from the worst observed value, pushing the
frame-free estimator below $1$~kcal/mol would require roughly
$6.6\times10^3$ shots per circuit, which on the measured rate card costs
about $\$23{,}400$ for this design; below $0.5$~kcal/mol, about
$2.7\times10^4$ shots and $\$94{,}000$. The productive direction is
therefore a lower-variance estimator, not a larger shot budget.

\subsection{Practical notes}

Three observations may help others budgeting work on this architecture.
Single-qubit gates were $58.9\%$ of this job's cost, so optimising
two-qubit count alone leaves most of the bill untouched. Device debiasing
multiplies cost by roughly $3.5\times$ and is on by default. Finally, we
encountered three defects in the provider's Qiskit integration: two in
multi-circuit result parsing and histogram scaling, since fixed upstream
but not at the time in a tagged release, and one in which a cost-estimate
field was read under the wrong key so that every estimate returned
\texttt{None}. Pipelines that budget from that field will silently see no
cost at all.

\section{Limitations}
\label{sec:limits}

We state these in the strongest form we can, because several of them bound
the claims above severely.

\textbf{This is not a blind variational calculation.} The
state-preparation angles are fit to the exact FCI Schmidt vectors, the
Schmidt coefficients used to assemble the energy are exact, and the
shared-frame reconstruction basis is spanned by the exact Schmidt
vectors and initialised at the exact frame. What is measured is whether a
noisy device plus mitigation can reproduce a \emph{known} state. This is a
meaningful test of a mitigation stack and a weaker claim than locating an
unknown ground state. This limitation was identified by an external
reviewer reading the source repository, and the frame-free estimator ---
which removes the reconstruction-side dependence though not the
circuit-construction one --- is reported as the headline for that reason.

\textbf{One molecule, one device, one submission.} The single-submission
condition that makes the pre-registration meaningful also prevents any
statement about run-to-run hardware reproducibility. The bootstrap
quantifies resampling uncertainty within one submission only.

\textbf{Chemical accuracy is not demonstrated.} The frame-free estimator
missed by $3.5\times$. The shared-frame estimator's point estimate falls
below the threshold but its bootstrap distribution is centred above it.

\textbf{The analytic correction contributes little to the shared-frame
estimator.} Disabling it entirely changes that estimator from $0.864$ to
$1.017$~kcal/mol. The joint frame fit, which is the component that depends
on the exact Schmidt vectors and which is initialised at the exact frame,
accounts for the remainder. The shared-frame number should therefore be
read as a measurement of what oracle-assisted reconstruction can recover,
not as a validation of the physics-based correction.

\textbf{The three reused circuits} were taken at $2000$ shots from an
earlier job on the same device, so the dataset is not uniform in shot
count.

\textbf{$K=6$ truncation.} The Schmidt rank is truncated; the resulting
classical floor is far below the errors discussed here but is not zero.

\section{Conclusion}

We pre-registered an analysis, ran one $81$-circuit submission on a
trapped-ion processor, and reported every estimator. A three-parameter
noise model fixed in advance predicted the raw energy error to $1\%$ and
the postselection retention to $0.4\%$. The correction built on that same
model then missed the oracle-free target by a factor of $4.16$, and the
goodness-of-fit statistic rose from $\approx2$ to $5.93$, flagging the
failure without reference to the known answer. Residual decomposition
localises the discrepancy to high-weight $XY$ strings in the largest
corrections and to a signed offset in one measurement group.

We did not reach chemical accuracy. We think the more useful output is the
pair of negative results: that raw-observable calibration does not certify
a correction model, and that at practical shot counts the oracle-free
estimator is limited by counting statistics rather than by device noise.
Both point away from more mitigation and towards lower-variance
estimators and honest fit reporting.

\begin{acknowledgments}
This work was supported by the IonQ Research Credits Program. We thank
Vadim Karpusenko (IonQ) for technical guidance throughout, for identifying
the native-gate compilation issue, for proposing the measurement-grouping
reduction, and in particular for the review of this repository's source
that identified the exact-state dependence described in
Sec.~\ref{sec:limits}. The pre-registration conditions for the hardware
run were set by IonQ and are reported verbatim.

In accordance with arXiv's policy on generative AI tools, we disclose that
a large language model (Anthropic Claude) was used as a coding and drafting
assistant during this work: for software implementation, for analysis
scripting, and for preparing this manuscript. All scientific claims,
experimental design decisions, pre-registration choices and
interpretations are the author's own, and the author takes full
responsibility for the content, including all numerical results, which were
generated by the committed code and are reproducible from the public
repository.
\end{acknowledgments}

\section*{Data and code availability}

All source code, the raw hardware counts for all $81$ circuits, the frozen
analysis script, and a complete iteration-by-iteration research log
including every negative result are available at
\url{https://github.com/Venkatallu11/quantum-chemistry-vqe}. The hardware
result is reproduced by
\texttt{task73\_hw\_analysis.py --counts task73\_hw\_counts.json
--design full}, and the residual decomposition by
\texttt{task82\_hw\_chi2\_residual\_breakdown.py}. Job identifier
\texttt{01a138e7-c5fd-711e-bd12-41e418290d36}.

\begin{thebibliography}{9}

\bibitem{peruzzo2014} A. Peruzzo, J. McClean, P. Shadbolt, M.-H. Yung,
X.-Q. Zhou, P. J. Love, A. Aspuru-Guzik, and J. L. O'Brien,
\emph{A variational eigenvalue solver on a photonic quantum processor},
Nat. Commun. \textbf{5}, 4213 (2014).

\bibitem{mcclean2016} J. R. McClean, J. Romero, R. Babbush, and
A. Aspuru-Guzik, \emph{The theory of variational hybrid
quantum-classical algorithms}, New J. Phys. \textbf{18}, 023023 (2016).

\bibitem{eddins2022} A. Eddins, M. Motta, T. P. Gujarati, S. Bravyi,
A. Mezzacapo, C. Hadfield, and S. Sheldon, \emph{Doubling the size of
quantum simulators by entanglement forging}, PRX Quantum \textbf{3},
010309 (2022).

\bibitem{pyscf} Q. Sun \emph{et al.}, \emph{PySCF: the Python-based
simulations of chemistry framework}, WIREs Comput. Mol. Sci. \textbf{8},
e1340 (2018).

\bibitem{temme2017} K. Temme, S. Bravyi, and J. M. Gambetta,
\emph{Error mitigation for short-depth quantum circuits},
Phys. Rev. Lett. \textbf{119}, 180509 (2017).

\bibitem{vandenberg2023} E. van den Berg, Z. K. Minev, A. Kandala, and
K. Temme, \emph{Probabilistic error cancellation with sparse
Pauli--Lindblad models on noisy quantum processors},
Nat. Phys. \textbf{19}, 1116 (2023).

\end{thebibliography}

\end{document}
